\section{Iteratives Verhältniswahlrecht mit garantiertem Gewinner}

\subsection{Verfahrensbeschreibung}

Das Verhältniswahlrecht tendiert dazu, Koalitionen zu erzwingen, weil keine Partei die absolute
Mehrheit erhält. Je mehr große Parteien es gibt, desto geringer ist die Wahrscheinlichkeit, dass
der Wähler eine klar bevorzugt. Das im Folgenden vorgeschlagene iterative Wahlverfahren sorgt
dafür, dass der Zwang zu zähneknirschenden Koalitionen mit faulen Kompromissen wegfällt. Es
garantiert, dass am Ende immer eine Regierungsmehrheit entsteht. Dabei gibt es aber auch den
Wählern, die im ersten Durchlauf \enquote{Verlierparteien} gewählt haben, in weiteren Durchläufen
die Chance, im Sinne des kleinsten Übels mitzubestimmen, wer am Ende regieren wird. Das
unterscheidet es von anderen mehrheitsbildenden Wahlverfahren.

\begin{enumerate}
  \item Im ersten Wahlschritt wird eine normale Verhältniswahl durchgeführt, bei der die Parteien
    nach der Anzahl ihrer prozentualen Stimmenanteile geordnet werden. Wenn eine Partei die
    absolute Mehrheit erhält oder sich eine Koalition bildet, die die absolute Mehrheit hat, ist
    das Verfahren beendet und die Partei oder Koalition mit der absoluten Mehrheit hat die Wahl
    gewonnen. Die \textbf{absolute Mehrheit} in diesem Sinne ist ein vorher definierter
    Prozentsatz, der leicht über 50\,\% liegt. Man wird ihn so festsetzen, dass die
    Regierungspartei/koalition ein paar Sitze mehr als die Hälfte bekommt, so dass sie
    rechnerisch auch noch die Mehrheit hat, wenn mal jemand krank ist.

  \item Wenn kein Gewinner feststeht, sind zwei Fälle zu unterscheiden:

    \begin{enumerate}
      \item Es standen nur zwei Parteien zur Auswahl. Dieser Fall unterteilt sich wieder in zwei
        Fälle:

        \begin{enumerate}[label=\roman*.]
          \item Es war der erste Wahlgang mit dieser Zusammensetzung. In diesem Fall wird
            Schritt~1 noch einmal wiederholt.

          \item Es ist schon der zweite ergebnislose Wahlgang mit nur zwei Parteien. In diesem
            Fall entscheidet ein Zufallsverfahren, beispielsweise minimale prozentuale
            Unterschiede im Ergebnis, oder das Los.
        \end{enumerate}

      \item Es standen mehr als zwei Parteien zur Auswahl. In diesem Fall werden die prozentualen
        Stimmanteile beginnend mit dem größten der Reihe nach aufsummiert, bis eine
        Zweidrittelmehrheit erreicht oder gerade überschritten ist. Das bedeutet, dass mindestens
        zwei Drittel der Wähler wollen, dass zumindest eine dieser Parteien regiert. Jetzt wird
        wieder Schritt~1 ausgeführt, aber nur mit den so bestimmten Parteien. Alle Wähler dürfen
        jetzt noch einmal über diese Parteien abstimmen.
    \end{enumerate}
\end{enumerate}

Das Verfahren wird so lange wiederholt, bis es einen Gewinner gibt. Das Verfahren terminiert
immer, da sich die Anzahl der zur Wahl stehenden Parteien in jedem Iterationsschritt reduziert,
bis gegebenenfalls nur noch zwei Parteien übrig sind. Der Gedanke hinter dem Verfahren ist der,
dass Menschen meistens auch eine Entscheidung treffen, wenn sie gezwungen werden, zwischen zwei
Alternativen zu entscheiden -- und wenn es nur die des kleinsten Übels ist. Bei vielen
Alternativen hingegen wählt der eine dies, der andere das. Koalitionen, wenn sie von Herzen
kommen, sind mehreren Wahlgängen vorzuziehen, da sie dem Wählerwillen am nächsten kommen. Es
besteht aber nicht mehr der Zwang, um jeden Preis zu koalieren, nur damit einer regieren kann.
Die Parteien können sich überlegen, wie ihre Chancen in einem erneuten Wahlgang stehen würden
und gegebenenfalls einen Koalitionspartner suchen oder nicht. Es kann durchaus vorkommen, dass
sich zwei Parteien inhaltlich nahestehen, aber eben zwei verschiedene Mannschaften sind, die um
die Gunst der Wähler eifern. In diesem Fall wäre eine Koalition auch aus Wählersicht anzuraten.

Wenn beim ersten Durchlauf bereits eine Partei oder Koalition eine Regierungsmehrheit ($\geq$
absolute Mehrheit) erzielen konnte, so erhält sie eine entsprechende Anzahl Sitze im Parlament
zugeteilt. Andernfalls bekommt der Gewinner aus mehreren Wahlgängen eine der absoluten Mehrheit
entsprechende Anzahl der Parlamentssitze zugesprochen. Die verbleibenden Sitze werden nach dem
ursprünglichen Verhältnis beim ersten Wahlgang zwischen den anderen Parteien aufgeteilt. So ist
sichergestellt, dass es immer eine Regierungsmehrheit gibt, aber auch kleinere Parteien in der
Opposition repräsentiert sind, solange ihr Stimmenergebnis für wenigstens einen Abgeordneten
ausreicht.

\underline{Beispiel:}

\bigskip
\noindent
\begin{minipage}[t]{0.48\textwidth}
\centering
1.~Wahlgang

\begin{tikzpicture}[x=2.5cm, y=0.08cm]
  \fill[diagrammHG] (0,0) rectangle (1,100);
  % Segmente von unten nach oben
  \fill[parteiA] (0,0)     rectangle (1,30);
  \fill[parteiB] (0,30)    rectangle (1,52.5);
  \fill[parteiC] (0,52.5)  rectangle (1,72.5);
  \fill[parteiD] (0,72.5)  rectangle (1,82.5);
  \fill[parteiE] (0,82.5)  rectangle (1,90.75);
  \fill[parteiF] (0,90.75) rectangle (1,96.75);
  \fill[parteiS] (0,96.75) rectangle (1,100);
  % Rahmen
  \draw (0,0) rectangle (1,100);
  % Prozentbeschriftungen
  \node[right] at (1,15)    {\small 30\,\%};
  \node[right] at (1,41.25) {\small 22{,}5\,\%};
  \node[right] at (1,62.5)  {\small 20\,\%};
  \node[right] at (1,77.5)  {\small 10\,\%};
  \node[right] at (1,86.6)  {\small 8{,}25\,\%};
  \node[right] at (1,93.75) {\small 6\,\%};
  % Parteikürzel links
  \node[left] at (0,15)    {A};
  \node[left] at (0,41.25) {B};
  \node[left] at (0,62.5)  {C};
  \node[left] at (0,77.5)  {D};
  \node[left] at (0,86.6)  {E};
  \node[left] at (0,93.75) {F};
  \node[left] at (0,98.4)  {S};
  % 66%-Linie
  \draw[thick] (-0.3,66) -- (1,66);
  \node[left] at (-0.3,66) {\small 66\,\%};
\end{tikzpicture}
\end{minipage}
\hfill
\begin{minipage}[t]{0.48\textwidth}
\centering
2.~Wahlgang

\begin{tikzpicture}[x=2.5cm, y=0.08cm]
  \fill[diagrammHG] (0,0) rectangle (1,100);
  % Segmente: A=45%, C=30%, B=25%
  \fill[parteiA] (0,0)  rectangle (1,45);
  \fill[parteiC] (0,45) rectangle (1,75);
  \fill[parteiB] (0,75) rectangle (1,100);
  \draw (0,0) rectangle (1,100);
  % Prozentbeschriftungen
  \node[right] at (1,22.5) {\small 45\,\%};
  \node[right] at (1,60)   {\small 30\,\%};
  \node[right] at (1,87.5) {\small 25\,\%};
  % Parteikürzel
  \node[left] at (0,22.5) {A};
  \node[left] at (0,60)   {C};
  \node[left] at (0,87.5) {B};
  % 66%-Linie
  \draw[thick] (-0.3,66) -- (1,66);
  \node[left] at (-0.3,66) {\small 66\,\%};
\end{tikzpicture}
\end{minipage}

\bigskip
\noindent
\begin{minipage}[t]{0.48\textwidth}
\centering
3.~Wahlgang

\begin{tikzpicture}[x=2.5cm, y=0.08cm]
  \fill[diagrammHG] (0,0) rectangle (1,100);
  % A=47.5%, C=52.5%
  \fill[parteiA] (0,0)    rectangle (1,47.5);
  \fill[parteiC] (0,47.5) rectangle (1,100);
  \draw (0,0) rectangle (1,100);
  \node[right] at (1,23.75) {\small 47{,}5\,\%};
  \node[right] at (1,73.75) {\small 52{,}5\,\%};
  \node[left] at (0,23.75) {A};
  \node[left] at (0,73.75) {C};
  \draw[thick] (-0.3,66) -- (1,66);
  \node[left] at (-0.3,66) {\small 66\,\%};
\end{tikzpicture}
\end{minipage}
\hfill
\begin{minipage}[t]{0.48\textwidth}
\centering
Ergebnisverteilung

\begin{tikzpicture}[x=2.5cm, y=0.08cm]
  \fill[diagrammHG] (0,0) rectangle (1,100);
  % C=51%, A=18.3%, B=13.8%, D=6.1%, E=5.1%, F=3.7%
  \fill[parteiC] (0,0)    rectangle (1,51);
  \fill[parteiA] (0,51)   rectangle (1,69.3);
  \fill[parteiB] (0,69.3) rectangle (1,83.1);
  \fill[parteiD] (0,83.1) rectangle (1,89.2);
  \fill[parteiE] (0,89.2) rectangle (1,94.3);
  \fill[parteiF] (0,94.3) rectangle (1,98);
  \draw (0,0) rectangle (1,100);
  % Labels
  \node[right] at (1,25.5)  {\small 51\,\%};
  \node[right] at (1,60.15) {\small 18{,}3\,\%};
  \node[right] at (1,76.2)  {\small 13{,}8\,\%};
  \node[right] at (1,86.15) {\small 6{,}1\,\%};
  \node[right] at (1,91.75) {\small 5{,}1\,\%};
  \node[right] at (1,96.15) {\small 3{,}7\,\%};
  % Parteikürzel
  \node[left] at (0,25.5)  {C};
  \node[left] at (0,60.15) {A};
  \node[left] at (0,76.2)  {B};
  \node[left] at (0,86.15) {D};
  \node[left] at (0,91.75) {E};
  \node[left] at (0,96.15) {F};
\end{tikzpicture}
\end{minipage}

\bigskip

Im obigen Beispiel gewinnt Partei~C im 3.~Wahlgang die Wahl, obwohl A im ersten Wahlgang die
stärkste Kraft war. Wenn A und B nach dem ersten Wahlgang eine Koalition eingegangen wären, wäre
die Wahl damit beendet gewesen und A und B hätten gemeinsam regiert. Da sie aber unvereinbare
Positionen oder Personen hatten, haben sie sich entschieden, in den nächsten Wahlgang zu gehen.
Am zweiten Wahlgang nehmen nur so viele Parteien teil, wie gerade nötig ist, um über zwei Drittel
der Wählerstimmen zu kommen. In diesem Fall waren das A, B und C. Beim zweiten Wahlgang konnte A
seinen Vorsprung ausbauen und kommt nun auf 45\,\%. Interessanterweise haben C und B die
Reihenfolge gewechselt, da offenbar mehr von den Wählern der im ersten Wahlgang ausgeschiedenen
Parteien Partei~C zugetan waren als Partei~B.
Da immer noch keine Partei die absolute Mehrheit hat und Koalitionen ausgeschlossen werden,
kommt es im 3.~Wahlgang zur Stichwahl zwischen A und C. Interessanterweise verliert A bei der
Stichwahl seine bisherige Führungsrolle, da die meisten ehemaligen B-Wähler C wählen. C hat
somit die Wahl gewonnen und bekommt die absolute Mehrheit zugesprochen. Da es sich nur um eine
zugesprochene absolute Mehrheit aufgrund von Stichwahlen handelt, ist das genaue prozentuale
Ergebnis der Stichwahlen irrelevant: Der Sieger erhält nur die zugesprochene Anzahl Sitze. (In
diesem Beispiel 51\,\%) Diese Regel ist deswegen wichtig, weil die Stichwahlen auch noch mit
einem viel größeren Unterschied ausgehen könnten. Z.\,B.\ könnte A auch 70\,\% der Stimmen
erhalten. Alle anderen Parteien wären dann absolut unterrepräsentiert, was nicht dem eigentlichen
Volkswillen entspräche. Wenn im ersten Wahlgang eine Partei oder Koalition 70\,\% der Stimmen auf
sich vereinen würde, wäre das eine andere Aussage. Dann wollen wirklich 70\,\% der Bevölkerung
trotz der Vielzahl der anderen Angebote, dass diese Partei(en) regieren.

Deswegen dürfen beim ersten Wahlgang erhaltene absolute Mehrheiten unverändert bestehen bleiben.

\subsection{Verworfene Alternative}

Anstatt die absolute Mehrheit der Stimmen zuzusprechen, könnte man auch einfach der
Gewinnerpartei des Iterativen Verfahrens das Recht der Gesetzgebung zusprechen. Wozu überhaupt
noch abstimmen? Sie gewinnen doch sowieso jede Abstimmung, da sie die absolute Mehrheit haben.
Das führt aber auf die Frage, wozu man überhaupt Parlamente braucht. Wenn die Fraktionen sowieso
immer genauso abstimmen, wie es ihr Fraktionsvorsitzender sagt, bräuchte man tatsächlich keine
Parlamente mehr. Dann würde es ausreichen, einen Fraktionsvorsitzendenrat zu bilden, wo die
Fraktionsvorsitzenden mit einem bestimmten Stimmgewicht abstimmen würden. Wenn der
Fraktionsvorsitzende der Regierungspartei dann auch noch der Regierungschef wäre, wäre die
Autokratie fertig.
Aber so ist es ja nicht. Die Parlamentarier können oft in geheimer Wahl abstimmen und die
Fraktionsvorsitzenden müssen ihre Fraktionsmitglieder überzeugen und können ihnen nicht
befehlen. Bei manchen Fragen kann man sich nicht einmal auf die eigenen Parteimitglieder
verlassen, sondern muss Verbündete in anderen Fraktionen suchen. Das ist genauso gewünscht,
damit die Macht von Einzelnen nicht zu groß wird. Wichtige Fragen sollen nicht von Einzelnen
entschieden werden, sondern von einer Gruppe von Menschen. Man geht aber davon aus, dass Menschen
mit gleichen Grundüberzeugungen, wie man es von Mitgliedern einer Partei annehmen kann, sich
meistens einigen werden. Manche Fragen werden allerdings auch innerhalb einer Partei so
kontrovers diskutiert, dass eine Einigung nicht immer möglich ist. Je größer und heterogener eine
Fraktion ist, desto schwieriger wird im Allgemeinen eine Einigung sein. Im Hinblick auf eine
Machtbegrenzung macht es deshalb schon Sinn, als Ausgleich für die Machterhöhung durch
Iterationswahlgewinn einen Zwang zur Einigung von mehr Menschen einzubauen. Das spricht dafür,
der Regierungspartei tatsächlich mehr Mandate zuzusprechen und nicht die Regel aufzustellen, dass
alle Gesetzesvorlagen der Regierungspartei als angenommen gelten, sofern nur die ganze
Regierungsfraktion zustimmt -- unabhängig von ihrer Größe.

\subsection{Mathematische Fragen zur Iterativen Verhältniswahl}

\begin{enumerate}
  \item Terminiert das Verfahren immer?
  \item Muss man den Schnitt immer nach $\tfrac{2}{3}$ machen, oder käme auch ein anderes
    Verhältnis infrage?
\end{enumerate}

Es bezeichne

\begin{itemize}[label={}]
  \item $A$ -- die Gesamtzahl der in einem Wahlgang abgegebenen Stimmen
  \item $A_i$ -- die Anzahl der Stimmen, die für die $i$-te Partei in diesem Wahlgang abgegeben
    wurde, wobei die Parteien absteigend nach ihren Stimmen geordnet sein sollen. D.\,h.\ wenn
    Partei $X$ im Wahlgang $j$ die $i$-te Partei war, muss sie das im Wahlgang $j+1$ nicht
    notwendigerweise auch wieder sein.
\end{itemize}

Für die relativen Stimmanteile $a_i = \frac{A_i}{A}$ gilt in jedem Wahlgang:

\begin{align}
  \text{I.} \quad & \sum_{i=1}^{n} a_i = 1
    \quad \text{($n$ -- Anzahl teilnehmender Parteien in diesem Wahlgang)} \\[4pt]
  \text{II.} \quad & a_k \geq a_l \quad \text{für } k < l
\end{align}

Um die Parteien zu bestimmen, die in den nächsten Wahlgang kommen, machen wir bei der relativen
Schnittstimmenzahl $r$ ($0 < r < 1$) einen Schnitt, so dass gilt:

\begin{align}
  \text{III.} \quad & \sum_{i=1}^{s} a_i \geq r \\[4pt]
  \text{IV.}  \quad & \sum_{i=1}^{s-1} a_i < r
\end{align}

Damit das Verfahren immer terminiert und sinnvolle Ergebnisse liefert, muss in jedem Durchgang
außer dem letzten (d.\,h.\ $n = 2$) $s < n$ und $s > 1$ gelten. Für $n = 2$ gilt Regel 2a, die
sicherstellt, dass es einen Gewinner gibt. Für alle anderen Durchgänge $j$ gilt: $n_{j+1} = s_j$.
Also solange $s_j < n_j$ erfüllt ist, wird $n$ in jedem Durchgang kleiner bis irgendwann 2
erreicht ist.
Ohne Beschränkung der Allgemeinheit können wir annehmen, dass in jedem Wahlgang $\forall a_i \neq
0$ gilt. Sollte in der Praxis doch jemals eine Partei in einem Wahlgang gar keine Stimme erhalten,
nehmen wir diese aus der Liste und reduzieren $n$ entsprechend, so dass immer $a_n \neq 0$ gilt.

Um zu beweisen, dass immer $s < n$ gilt, nehmen wir zuerst das Gegenteil an, d.\,h.

\medskip
Annahme A1: $s = n$

\medskip\noindent
und versuchen diese Behauptung zum Widerspruch zu führen. Da nie $s > n$ sein kann, muss dann
$s < n$ gelten.

Aus A1 und IV folgt
\begin{equation}
  \text{V.} \quad \sum_{i=1}^{n-1} a_i < r
\end{equation}

Addition von $a_n$ auf beiden Seiten ergibt
\begin{equation}
  \text{V'.} \quad \sum_{i=1}^{n-1} a_i + a_n < r + a_n
\end{equation}

aus (I) $\Rightarrow$ $1 < r + a_n$ $\Rightarrow$ $1 - a_n < r$
\begin{equation}
  \text{VI.} \quad a_n > 1 - r
\end{equation}

aus (II) $\Rightarrow$ $a_i \geq a_n$ für $i < n$
\begin{equation}
  \text{VII.} \quad \sum_{i=1}^{n-1} a_i \geq \sum_{i=1}^{n-1} a_n
\end{equation}

(VI.) $\Rightarrow$
\begin{equation}
  \text{VIII.} \quad \sum_{i=1}^{n-1} a_n > \sum_{i=1}^{n-1}(1-r) = (n-1)(1-r)
\end{equation}

(VII, VIII, VI) $\Rightarrow$
\begin{equation}
  \text{IX.} \quad \sum_{i=1}^{n} a_i = \sum_{i=1}^{n-1} a_i + a_n
    > (n-1)(1-r) + a_n > (n-1)(1-r) + (1-r) = n(1-r)
\end{equation}

Also $\displaystyle\sum_{i=1}^{n} a_i > n(1-r)$.

Das ist ein Widerspruch zu (I) genau dann wenn
\begin{equation}
  \text{X.} \quad n(1-r) \geq 1
\end{equation}

weil $\displaystyle\sum_{i=1}^{n} a_i > x$ nur für $x < 1$ gelten kann, da
$\displaystyle\sum_{i=1}^{n} a_i = 1$ ist.

A1 ist also nicht immer falsch, sondern die Ungültigkeit hängt von $n$ und $r$ ab.
Um sicherzustellen, dass A1 falsch ist und damit $s < n$, muss gelten:
\begin{equation}
  \text{XI.} \quad n \geq \frac{1}{(1-r)} \quad \text{bzw.} \quad r \leq 1 - \frac{1}{n}
\end{equation}

Nach Voraussetzung ist $n > 2$. Für $n = 2$ haben wir unsere Sonderregelung, nach der wir den
Verfahrensgewinner bestimmen. Wir brauchen, dass A1 falsch ist für:
\begin{equation}
  \text{XII.} \quad n \geq \frac{1}{(1-r)} = 3
\end{equation}
\begin{equation}
  \text{XIII.} \quad r = \frac{2}{3}
\end{equation}

Das bedeutet, dass sich die Anzahl der zur Verfügung stehenden Parteien in jedem
Iterationsschritt reduziert, bis nur noch zwei Parteien übrig bleiben -- vorausgesetzt wir
schneiden bei zwei Dritteln. Wenn wir z.\,B.\ erst nach drei Vierteln schneiden würden, bräuchten
wir eine Sonderregelung schon für drei Parteien im Wahlgang und nicht erst für zwei, da eine
Reduktion auf zwei nicht mehr allgemeingültig garantiert werden könnte.

Man kann sich das auch anschaulich gut vorstellen:

Wenn drei Parteien im Spiel sind, die alle genau gleich viele Stimmen erhalten, dann können wir
das Verfahren nicht anwenden, da Voraussetzung~IV nicht erfüllt ist. Wenn die Parteien aber auch
nur minimale Stimmenunterschiede aufweisen (wovon bei Wahlen mit vielen Wählern und wenigen zur
Wahl Stehenden auszugehen ist), wird eine Partei weniger als ein Drittel der Stimmen erhalten und
die anderen beiden zusammen damit über zwei Drittel kommen, womit im nächsten Wahlgang nur die
beiden drinbleiben. Würden wir aber bei drei Vierteln erst den Schnitt machen, würden bei drei
Parteien minimale Ergebnisunterschiede keinesfalls eine Partei unter ein Viertel kommen lassen.
Da müssten sich die Ergebnisse schon deutlich unterscheiden, was sein kann, aber nicht sein muss.
Wären hingegen vier Parteien im Spiel, würden minimale Ergebnisunterschiede auch bei einem drei
Viertel Schnitt eine Partei rausfliegen lassen.

\subsection{Verfassungsänderungen}

Aus gutem Grund dürfen Verfassungsänderungen nur mit einer besonders großen Mehrheit, der
\enquote{Verfassungsmehrheit}, die in Deutschland zwei Drittel beträgt, durchgeführt werden.
Verfassungsänderungen sollen nur mit einem breiten gesellschaftlichen Konsens möglich sein --
auch über Parteigrenzen hinweg. Die Zwei-Drittel-Regelung kann nicht auf Parlamente mit
zugesprochener absoluter Mehrheit angewendet werden, da dies einen größeren gesellschaftlichen
Konsens vortäuschen würde, als tatsächlich vorhanden. Die Regierungspartei hat im Sinne des
kleinsten Übels mehr Macht bekommen, als ihr eigentlich zustünde, damit wenigstens einer etwas
entscheiden kann, auch wenn dadurch am Ende vielleicht nur die zweitbeste Lösung umgesetzt wird.
Wenigstens wird etwas umgesetzt, was aus einem Guss ist.
Dieses Prinzip des kleinsten Übels sollte aber nicht bei Grundsatzfragen zur Anwendung kommen.
Hier muss richtig entschieden werden, da die Folgen von Fehlentscheidungen zu gravierend wären.
Deswegen schlage ich vor, dass eine Verfassungskammer gebildet wird, die die Zusammensetzung hat,
wie sie im ersten Wahlgang vom Wähler festgelegt worden ist. Alle Verfassungsänderungen müssen
mit Zweidrittelmehrheit in dieser Kammer bestätigt werden. Wenn dazu größere überparteiliche
Absprachen notwendig sind, dann ist das so und durchaus im Sinne des Erfinders.
Die Verfassungskammer kann auch über Verfassungsänderungen hinaus genutzt werden, um andere
verfassungsrechtlich relevante Entscheidungen zu treffen, wie z.\,B.\ die Wahl der
Verfassungsrichter.
Bei allen verfassungsrelevanten Entscheidungen geht man davon aus, dass die Verfassung
eigentlich schon gut so ist, wie sie ist. Deswegen macht man es bewusst schwer, sie zu ändern,
damit man sie nicht aus Versehen verschlimmbessert. Sie ist einmal durch einen breiten
gesellschaftlichen Konsens zustande gekommen\footnote{Der Konsens, der zur Bildung der Verfassung
der Bundesrepublik Deutschland geführt hat, mag nicht freiwillig zustande gekommen sein, sondern
durch einen verlorenen Krieg. Er ist aber doch im Nachhinein von Millionen von Bundesdeutschen im
Grundsatz akzeptiert worden, so dass man schon von einem gesamtgesellschaftlichen Konsens
sprechen kann.} und nur ein mindestens genauso breiter gesellschaftlicher Konsens sollte sie
wieder ändern können. Wenn der nicht zustande kommt, dann muss sie so bleiben, wie sie ist.
Sollte ein gleich großer Konsens über Jahrzehnte nicht zustande kommen, dann ist es eben so und
wir leben weiter in dem von den (geistigen) Vätern vorgegebenen Rahmen. Das Prinzip stößt nur an
Grenzen bei Verfassungsrichtern, die im Gegensatz zu niedergeschriebenen Verfassungen mit der
Zeit altern und ausgetauscht werden müssen. Deswegen muss auch da ein Verfahren geben, wie neue
Verfassungsrichter bestimmt werden können, die in die Fußstapfen ihrer Vorgänger treten, ohne
dass notwendigerweise die geforderte Zweidrittelmehrheit zustande kommt. Das Naheliegendste ist
tatsächlich, die amtierenden Verfassungsrichter sich selbst Nachfolger suchen zu lassen. Sie
sollen sie nicht nur vorschlagen können, sondern sie selbst bestimmen, außer die Politik greift
mit Zweidrittelmehrheit ein. Von sich aus werden sie Nachfolger einsetzen, die in ihrem Geiste
entscheiden werden.
